\section{Confluencia de acción, área, información y gravitación}
\label{sec:confluencia}

La escala areal de esta sección es la misma que determina la
composición longitudinal y radial de \ref{g26:sec:rigidez-radial-es}:
\begin{equation}
 \ell_P^2=L^2=\frac{G_a\hbar_a}{c_{\rm int}^3},\qquad
 E_{P,a}=\frac{\hbar_a c_{\rm int}}{L}
        =k_B\Theta_{{\rm clk},a}\log3.
 \label{eq:ii-g26-longitud-energia-area}
\end{equation}
En la sección de retorno, \(E_{P,\rm ret}=E_L\). El radio variable
\(R\) del horizonte conserva su significado propio; la condición
\(R=L\) especifica la celda particular estudiada a continuación.
El incremento \(\delta\mathcal A=\gamma_{\HMT}a_0L^2\) transmite
esa misma unidad areal a la incidencia, manteniendo la acción, el
transductor térmico y las ocupaciones del estado reducido.

\subsection{Energía y entropía de la celda de media acción}

Fijamos una misma sección orientada de acción y sus secciones
gravitatoria y térmica correspondientes. La reciprocidad anterior permite expresar la misma celda mediante
acción, duración y longitud. Con
\(\ell_P^2=\hbar G/c^3\), \(t_P=\ell_P/c\) y
\(E_P=\hbar/t_P\), se tiene
\[
 \frac{c^4\ell_P}{G}=\frac{\hbar}{t_P}=E_P.
\]
La igualdad sigue de sustituir la escala areal ya construida. En la
realización de horizonte definida por \eqref{eq:ii-horizonte-normalizado}, las cantidades
se normalizan como
\begin{equation}
 E_\Lambda(R)=\frac{c^4R}{2G},\qquad
 S_H(R)=\frac{k_B\pi R^2}{\ell_P^2},\qquad
 T_H(R)=\frac{\hbar c}{2\pi k_BR}.
 \label{eq:ii-horizonte-normalizado}
\end{equation}
La temperatura corresponde a la normalización de horizonte cosmológico
de esta realización. El factor \(2\pi\) y la definición del radio
forman parte del conjunto de ecuaciones que se utiliza.

\begin{theorem}[Identidad de la celda gravitatoria de media acción]
En la realización \eqref{eq:ii-horizonte-normalizado},
\begin{align}
 T_H(R)S_H(R)&=E_\Lambda(R),\\
 T_H(\ell_P)S_H(\ell_P)&=\frac{E_P}{2}
 =\frac{k_B\Theta_{\rm clk}\log3}{2},\\
 E_\Lambda(\ell_P)t_P&=\frac{\hbar}{2},\qquad
 E_\Lambda(\ell_P)T_{108}=\frac h2.
\end{align}
\end{theorem}
\begin{proof}
Multiplicar la temperatura por la entropía cancela la normalización
térmica y da \(\hbar cR/(2\ell_P^2)=c^4R/(2G)\).
En \(R=\ell_P\), la expresión es \(\hbar c/(2\ell_P)
=\hbar/(2t_P)=E_P/2\). La identidad térmica del ciclo da
el segundo miembro. Finalmente, \(T_{108}=2\pi t_P\) y
\(h=2\pi\hbar\) producen las dos relaciones de acción.
\end{proof}

El teorema reúne cantidades que previamente pertenecían a descripciones
distintas. La energía de la celda, el producto de temperatura y entropía,
la información trítica y la acción integrada sobre su duración son
lecturas compatibles bajo la normalización declarada. La constante
gravitatoria interviene en la escala areal; la acción fija la energía
del período; Boltzmann transporta la información adimensional a la
recta de entropía. La demostración muestra cómo estas funciones se
articulan en una misma identidad.

\subsection{Conversión entre energía de decisión e incremento areal}
\label{sec:ii-conversion-informacion-area}

La energía por decisión trítica y el incremento elemental de área se
obtienen mediante lectores distintos del desarrollo anterior:
\begin{equation}
 E_P=k_B\Theta_{\rm clk}\log3=\frac{\hbar}{t_P},
 \qquad
 \delta\mathcal A=\gamma_{\HMT}a_0\ell_P^2.
 \label{eq:ii-confluencia-cuanto-energia-area}
\end{equation}
Se mantienen la misma sección orientada y la rama positiva
\(\hbar,G,\gamma_{\HMT},a_0,t_P>0\). La primera cantidad
pertenece a la recta de energía y la segunda a la recta de área.
Su cociente define el coeficiente de conversión
\begin{equation}
 \mathcal T_{I/A}
 :=\frac{E_P}{\delta\mathcal A}
 =\frac{k_B\Theta_{\rm clk}\log3}
        {\gamma_{\HMT}a_0\ell_P^2},
 \qquad [\mathcal T_{I/A}]=[\text{energía}]/[\text{área}].
 \label{eq:ii-confluencia-tension}
\end{equation}
Esta definición explicita la normalización necesaria para transportar
una energía de decisión a una resolución geométrica. Todos sus factores
están especificados antes de formar el cociente.

\begin{proposition}[Expresión gravitatoria del coeficiente de conversión]
\label{prop:ii-confluencia-tension}
En la misma realización dimensional,
\begin{equation}
 \mathcal T_{I/A}
 =\frac{c^3}{\gamma_{\HMT}a_0G t_P}
 =\frac{c^4}{\gamma_{\HMT}a_0G\ell_P}.
 \label{eq:ii-confluencia-tension-gravedad}
\end{equation}
\end{proposition}
\begin{proof}
Sustituir \(E_P=\hbar/t_P\) y \(\ell_P^2=\hbar G/c^3\)
en \eqref{eq:ii-confluencia-tension} cancela la sección de acción
común y da el primer miembro. La identidad \(\ell_P=ct_P\)
produce el segundo. Se conserva la dependencia de \(t_P\) y
\(\ell_P\) respecto de la realización previamente construida.
\end{proof}

En el corte mínimo de \eqref{eq:ii-area-corte-informacion}, la
realización de una decisión uniforme por enlace da
\begin{equation}
 S_{\rm corte}^{\rm fis}=81k_B\log3,
 \qquad
 E_{\rm corte}^{(1)}=81E_P
 =\Theta_{\rm clk}S_{\rm corte}^{\rm fis}.
 \label{eq:ii-confluencia-corte-energia}
\end{equation}
La primera igualdad utiliza la capacidad demostrada y el estado producto
uniforme; la segunda asigna la energía de decisión del mismo ciclo a
cada uno de sus ochenta y un enlaces. Multiplicar la identidad térmica
por ochenta y uno prueba la ecuación de estado de esta realización.
El estado sectorial de treinta y seis enlaces conserva su propia
función de partición, su entropía y las variaciones areales de
\eqref{eq:ii-area-ley-finita}.

La composición identifica las funciones complementarias de las
constantes: la acción fija energía por duración; el carácter angular
y Barbero--Immirzi fijan la resolución sectorial de área; Boltzmann
realiza dimensionalmente la información; el producto \(G\hbar\)
determina la unidad areal. El coeficiente \(\mathcal T_{I/A}\)
conserva estas dependencias al expresar su relación en unidades de
energía por área.

\subsection{Significado físico de la confluencia areal}

En el horizonte de radio \(R\), el área es \(A_H=4\pi R^2\).
La misma relación entrópica admite las dos expresiones
\[
 \frac{S_H}{k_B}=\frac{A_H}{4\ell_P^2}
 =\frac{\pi R^2c^3}{G\hbar}.
\]
El primer cociente compara el área macroscópica con la unidad areal
de la construcción. El segundo explicita la composición de esa unidad
mediante acción y gravitación. Al transportar las secciones de modo
que \(G_a\hbar_a\) permanezca fijo, la lectura geométrica conserva
esa normalización. La entropía adimensional del mismo horizonte queda
invariante; sus expresiones dimensionales conservan las bases que
utilizan.

El parámetro de Barbero--Immirzi actúa en una resolución diferente
de ese mismo desarrollo: la incidencia orientada organiza el espectro
elemental de área y el estado reducido describe la información
accesible en el corte. Las hojas areal y volumétrica aportan el
transporte geométrico y la medida de historias descritos en
\cref{sec:ii-hojas-av}. La comparación entre un espectro microscópico
y una entropía de horizonte conserva como datos su mapa de realización,
el estado y la escala. Estas precisiones permiten interpretar la
relación física en cada nivel y evitan sustituir una correspondencia
de observables por una identificación indiscriminada de sus espacios.

La unidad teórica del argumento reside en funciones complementarias:

La unidad areal reúne el producto $G\hbar$; la incidencia orientada
organiza su resolución elemental mediante $\gamma_{\HMT}a_0$; el
estado y la bipartición determinan la información accesible al
subsistema. Las ecuaciones de respuesta
\eqref{eq:ii-respuesta-entropica}--\eqref{eq:ii-respuesta-areal}
explicitan esa articulación en la colección sectorial. La inversión
\eqref{eq:ii-inversion-entropia} conserva la entropía y transforma
el área según~\eqref{eq:ii-area-complementaria}: las dos lecturas
distinguen propiedades diferentes de la misma configuración.

En esta articulación,
la dinámica produce y conserva la historia; el lector determina qué
información es accesible; la acción relaciona evolución y energía;
Boltzmann realiza su expresión térmica; la gravitación interviene
en la escala geométrica con la que se mide el área. La identidad
\(T_HS_H=E_\Lambda\) muestra una confluencia exacta de esas funciones
en la realización declarada. La teoría física se formula así desde
los objetos matemáticos y desde sus transformaciones efectivas.

\subsection{Entropía espectral y conservación bajo transporte}

La entropía de un estado reducido depende de su espectro y de la
bipartición que define la reducción. Para precisar la continuación
hacia dualidad, sean \(\mathcal H_A,\mathcal H_B\) los dos factores
y \(\rho_{AB}\) un estado con entropías finitas. Un transporte
factorizado \(U_A\otimes U_B\) produce
\[
 \rho'_A=U_A\rho_AU_A^*,\qquad S(\rho'_A)=S(\rho_A).
\]
La primera identidad sigue de la definición de traza parcial y la
segunda de la invariancia del espectro. Para un transporte unitario
general, la formulación equivalente consiste en transportar también
la inclusión del álgebra de observables del subsistema. Esta precisión
identifica el dato geométrico que debe acompañar a una afirmación
de conservación de entrelazamiento.

La reciprocidad de la elipse y el intercambio momento--enrollamiento
conservan las cantidades indicadas en sus respectivos teoremas. El
transporte de la bipartición determina cómo utilizar esas identidades
en una lectura de entropía. De este modo, la continuación modular
mantiene el contenido informacional del artículo y conserva separados
el espectro energético, la orientación de la acción y la reducción
algebraica que define el entrelazamiento.

\subsection{Multiplicidad efectiva y realización hiperbólica}

En la celda anterior, \(S_H(\ell_P)/k_B=\pi\). La multiplicidad
efectiva asociada a esa entropía es
\[
 \Omega_{\rm eff}=\exp(S_H/k_B)=e^\pi.
\]
Si \(H\) es la involución hiperbólica del TRIT, con autovalores
\(+1,-1\), la exponencial \(\exp(\pi H)\) tiene espectro
\(\{e^\pi,e^{-\pi}\}\). La coincidencia del autovalor expansivo
con la multiplicidad efectiva sitúa una relación espectral específica.
Su significado microscópico conserva como dato adicional el mapa entre
el espacio del horizonte y la representación hiperbólica. Esta
localización permite desarrollar la correspondencia sin confundir
el valor espectral con la totalidad de su espacio de realización.
